%==============================================================================
% A1 POSTER — Bayesian spatio-temporal models for Italian municipal migration
% Department of Statistical Sciences, University of Padua
%==============================================================================

\documentclass[final]{beamer}

\usepackage[
  orientation=portrait,
  size=a1,
  scale=1.18
]{beamerposter}

%------------------------------------------------------------------------------
% Packages
%------------------------------------------------------------------------------
\usepackage[utf8]{inputenc}
\usepackage[T1]{fontenc}
\usepackage{amsmath,amssymb,mathtools,bm}
\usepackage{booktabs}
\usepackage{array}
\usepackage{graphicx}
\usepackage{multicol}
\usepackage{xcolor}
\usepackage{ragged2e}
\newcommand{\gap}{\vspace{1em}}
\usepackage{tikz}
\usepackage{qrcode}
\usepackage{hyperref}

%------------------------------------------------------------------------------
% Geometry / spacing
%------------------------------------------------------------------------------
\setlength{\paperwidth}{59.4cm}
\setlength{\paperheight}{84.1cm}

\setbeamersize{
  text margin left=1.4cm,
  text margin right=1.4cm
}

\setlength{\columnsep}{0.95cm}
\setlength{\columnseprule}{0pt}

%------------------------------------------------------------------------------
% University of Padua inspired palette
%------------------------------------------------------------------------------
\definecolor{padua}{RGB}{155,0,20}
\definecolor{paduaDark}{RGB}{110,0,15}
\definecolor{lightgray}{RGB}{246,246,248}
\definecolor{midgray}{RGB}{90,90,90}

\setbeamercolor{background canvas}{bg=white}
\setbeamercolor{normal text}{fg=black,bg=white}
\setbeamercolor{structure}{fg=padua}
\setbeamercolor{alerted text}{fg=padua}

\setbeamercolor{block title}{fg=white,bg=padua}
\setbeamercolor{block body}{fg=black,bg=lightgray}

\setbeamercolor{block title alerted}{fg=white,bg=paduaDark}
\setbeamercolor{block body alerted}{fg=black,bg=white}

\setbeamertemplate{navigation symbols}{}

%------------------------------------------------------------------------------
% Fonts
%------------------------------------------------------------------------------
\setbeamerfont{title}{size=\huge,series=\bfseries}
\setbeamerfont{author}{size=\Large}
\setbeamerfont{institute}{size=\large}
\setbeamerfont{block title}{size=\Large,series=\bfseries}
\setbeamerfont{block body}{size=\normalsize}
\setbeamerfont{caption}{size=\footnotesize}

%------------------------------------------------------------------------------
% Short helper commands
%------------------------------------------------------------------------------
\newcommand{\model}[1]{\textbf{#1}}
\newcommand{\smallgap}{\vspace{0.25cm}}
\newcommand{\microgap}{\vspace{0.12cm}}

%------------------------------------------------------------------------------
% Poster metadata
%------------------------------------------------------------------------------
\title{Bayesian spatio-temporal modeling of net migration across Italian municipalities}

\author{C. Franceschini \and S.G. Parolin \and T. Secco}

%\institute{
%Department of Statistical Sciences, University of Padua
%}

\date{}

%==============================================================================

\begin{document}
\begin{frame}[t]

%------------------------------------------------------------------------------
% LOGO + QR CODE -- overlay, non occupano spazio
%------------------------------------------------------------------------------
\begin{tikzpicture}[remember picture,overlay]

  % Logo dipartimento -- alto a destra
  \node[
    anchor=north east,
    xshift=-0.65cm,
    yshift=-0.5cm
  ] at (current page.north east)
  {
    \includegraphics[width=0.15\textwidth]
    {Dipartimento_Scienze_statistiche.png}
  };
% QR code -- alto a sinistra
\node[
  anchor=north west,
  xshift=0.8cm,
  yshift=-0.75cm,
  draw=padua,
  line width=1.5pt,
  rounded corners=3pt,
  inner sep=0.12cm
] (qr) at (current page.north west)
{
  \includegraphics[width=0.08\textwidth]
  {qrcode.jpeg}
};

% Titolo QR
\node[
  anchor=south,
  yshift=-0.12cm,
  text=padua
] at (qr.north)
{
  {\small\bfseries Learn more}
};
\end{tikzpicture}

%------------------------------------------------------------------------------
% TITLE AREA
%------------------------------------------------------------------------------
\begin{center}
\vspace{-1cm}
{\fontsize{45}{45}\selectfont\bfseries\color{padua}
Bayesian spatio-temporal modeling of \\ net migration across Italian municipalities\par}

\vspace{0.5cm}

{\fontsize{28}{30}\selectfont\bfseries\color{paduaDark}
Associations, spatio-temporal structure, and territorial heterogeneity\par}

\vspace{0.5cm}

{\fontsize{20}{22}\selectfont
C. Franceschini
\qquad
S.G. Parolin
\qquad
T. Secco\par}

\end{center}

\vspace{0.10cm}
\hrule height 1.5pt
\vspace{0.25cm}


%==============================================================================
% THREE-COLUMN LAYOUT
%==============================================================================

\begin{columns}[t,totalwidth=\textwidth]

%==============================================================================
% COLUMN 1: EMPIRICAL SETTING, PREDICTORS & VALIDATION STRATEGY
%==============================================================================
\begin{column}{0.324\textwidth}

%------------------------------------------------------------------------------
\begin{block}{1. Italian municipal net migration}

\justifying
We investigate the spatio-temporal dynamics of \textbf{municipal net migration in Italy} using a comprehensive nationwide panel spanning 16 consecutive years:
\[
S = 7{,}896 \text{ municipalities}, \qquad T = 16 \text{ years (2004--2019)},
\]
\[
N = 126{,}336 \text{ municipality--year observations}.
\]
The response is the municipal \textbf{net migration rate per $1{,}000$ residents},
globally standardized before model fitting:
\[
y_{st} = 1{,}000 \times \frac{\text{Net Migration}_{st}}{\text{Resident Population}_{st}}.
\]

%This demographic normalization prevents high-density metropolitan areas from dominating the variance structure purely due to population scale.

\vspace{0.15cm}
\begin{center}
  \includegraphics[width=0.95\linewidth]{Poster_figures/raw_migration_map_comuni_en.png}
  \vspace{0.15cm}
  \parbox{0.96\linewidth}{\footnotesize \textit{Figure 1: Observed municipal net migration rate per $1{,}000$ residents across Italian municipalities in 2004, 2007, 2012, and 2017. %Highlights structural North--South disparities, depopulation of inland mountain belts, and the nationwide contraction following the 2008 crisis.
  }}
\end{center}

\gap
\justifying
\textbf{Covariates ($\mathbf{x}_{st}$):} 10 %standardized 
covariates covering 
\begin{itemize}
    \item \textbf{labor-market conditions:}
    employment and unemployment rates;

\item \textbf{demographic and economic dynamics:}
mean age, change in the foreign-resident share, log-differences in income, population density, and GDP;

    \item \textbf{territorial inequality and accessibility:}
    income inequality, landslide exposure, access to essential services.
\end{itemize}
\end{block}

\smallgap

%------------------------------------------------------------------------------
\begin{block}{2. Research questions}

\justifying
\begin{itemize}

\item How are \textbf{demographic, economic, and territorial characteristics} associated with Italian municipal net migration?

\item What \textbf{local spatial variation and temporal persistence} emerge in net migration at the italian municipal scale?

\item Does allowing for \textbf{local heterogeneity} reveal spatially varying covariate associations \textbf{beyond a common national specification}?

\end{itemize}

\end{block}
\smallgap
%------------------------------------------------------------------------------

\begin{block}{3. Areal topology and Leroux spatial precision}

\justifying
Municipalities form a \textbf{highly granular and irregular areal system}. %: nearby areas may share similar migration dynamics, while each municipality is repeatedly \textbf{observed over time}. 
\textbf{Spatial dependence} is represented through a binary \textbf{contiguity graph}, on harmonized 2024 municipal borders:

\[
W_{ij}=
\begin{cases}
1, & i \text{ and } j \text{ share a boundary},\\
0, & \text{otherwise}.
\end{cases}
\]

The graph contains $|V|=7{,}896$ municipalities and
$|E|=22{,}609$ undirected edges. \\


%The network comprises $|V| = 7{,}896$ vertices and $|E| = 45{,}218$ shared borders.\\

\noindent We adopt the \textbf{Leroux CAR specification} for the spatial precision matrix:
\[
Q_S(\rho)=\rho(D-W)+(1-\rho)I_S,
\qquad \rho\in(0,1).
\] 
with 
%nota: Nel modello finale abbiamo \rho \sim \mathrm{Beta}(6,1) e una beta ha supporto (0,1)
\[
\qquad D_{ii} = \sum_j W_{ij}, \qquad \bar{d} = 5.73 \text{ neighbors}.
\]

%No artificial connections are introduced for geographically isolated municipalities.
\textbf{Key properties:}
\begin{itemize}

 \item as $\rho\rightarrow0$, the spatial effects approach independence,
  while increasing $\rho$ strengthens spatial smoothing;

  \item for $\rho\in(0,1)$, $Q_S(\rho)$ is positive definite and defines a
  proper CAR prior;

  \item as $\rho\rightarrow1$, the model approaches the intrinsic CAR
  specification, for which $D-W$ is singular;

  \item 14 municipalities have no terrestrial neighbors and are retained as
  isolated components, with no artificial marine links.

\end{itemize}

\textbf{Temporal dependence} is introduced through an \textbf{autoregressive evolution} of the spatial random effects.

\smallgap
\end{block}

%------------------------------------------------------------------------------
\begin{block}{References}

\scriptsize
\justifying
\begin{multicols}{2}

\textbf{Mozdzen et al.} (2022).
Bayesian modeling and clustering for spatio-temporal areal data:
An application to Italian unemployment.
\textit{Spatial Statistics}, 52, 100715.

\textbf{Orozco-Acosta, Adin \& Ugarte} (2021).
Scalable Bayesian modeling for smoothing disease risks in large spatial
data sets using INLA.
\textit{Spatial Statistics}, 41, 100496.

\textbf{Orozco-Acosta, Adin \& Ugarte} (2023).
Big problems in spatio-temporal disease mapping: methods and software.
\textit{Computer Methods and Programs in Biomedicine}, 231, 107403.

\textbf{Rue, Martino \& Chopin} (2009).
Approximate Bayesian inference for latent Gaussian models by using
integrated nested Laplace approximations.
\textit{JRSS-B}, 71(2), 319--392.

\end{multicols}

\end{block}

\end{column}

%==============================================================================
% COLUMN 2: MODELING BACKBONE & COMPARATIVE EVALUATION
%==============================================================================
\begin{column}{0.324\textwidth}
%------------------------------------------------------------------------------
\begin{block}{4. The global spatio-temporal backbone}

The global homogeneous model decomposes the standardized response into
common covariate effects, a \textbf{national temporal trend}, a
spatio-temporal latent field, and Gaussian noise:
\[
y_{st}^*
=
\alpha
+
\mathbf{x}_{st}^{*\prime}\bm{\beta}
+
f_t
+
w_{st}
+
\epsilon_{st},
\qquad
\epsilon_{st}\sim\mathcal{N}(0,\sigma_\epsilon^2).
\]

The \textbf{common temporal component} follows a first-order random walk:
\[
f_t-f_{t-1}
\sim
\mathcal{N}(0,\tau_{\mathrm{RW1}}^2),
\qquad t=2,\ldots,T,
\qquad
\sum_{t=1}^{T}f_t=0.
\]

\microgap
\textbf{Spatio-temporal latent field $\mathbf{w}_t$:}
\[
\mathbf{w}_1
\sim
\mathcal{N}\!\left(
\mathbf{0},
\tau_{ST}^2Q_S(\rho)^{-1}
\right),
\]
\[
\mathbf{w}_t\mid\mathbf{w}_{t-1}
\sim
\mathcal{N}\!\left(
\xi\mathbf{w}_{t-1},
\tau_{ST}^2Q_S(\rho)^{-1}
\right),
\qquad t=2,\ldots,T.
\]

This yields the separable Kronecker-structured joint precision for
$\mathbf{w}\in\mathbb{R}^{ST}$:
\[
Q_{ST}(\rho,\xi,\tau_{ST}^2)
=
\frac{1}{\tau_{ST}^2}
Q_T(\xi)\otimes Q_S(\rho),
\]
where $Q_T(\xi)$ is the sparse tridiagonal precision induced by the first-order autoregressive evolution.

%The $IT=126{,}336$-dimensional latent field is handled through sparse precision-matrix factorizations.

\vspace{0.5cm}
\textbf{Two matched computational implementations:}
\begin{itemize}
  \item \model{M1}: Integrated Nested Laplace Approximation (INLA);
  \item \model{M2}: custom Metropolis-within-Gibbs sampler adapted from
  Mozdzen et al.\ (2022) (20k iter., 10k burn-in, thin 10).
\end{itemize}

\vspace{0.5cm}
Following Mozdzen et al.\ (2022), fixed effects have
$\mathcal{N}(0,1)$ priors.
$\sigma_\epsilon^2,\tau_{ST}^2\sim\mathrm{IG}(3,2)$;
$\rho\sim\mathrm{Beta}(6,1)$;
$\xi\sim\mathrm{Unif}(-1,1)$.
\end{block}

\vspace{0.5cm}

%------------------------------------------------------------------------------
\begin{block}{5. Alternative modeling strategies}

\justifying
%We benchmark the global backbone against three alternatives:

\begin{itemize}
\item \textbf{Partitioned model (INLA) (\model{M3} Overlap, \model{M4} Disjoint;
Orozco-Acosta et al., 2023):}\\

\noindent 107 province-specific models allow fixed effects and spatio-temporal
parameters to vary locally. \model{M3} adds first-order neighboring
municipalities to reduce boundary effects; \model{M4} uses disjoint provinces.

\item \textbf{Bayesian nonparametric clustering (\model{M5 BSTC-DP}; Mozdzen et al., 2022):}\\

\noindent municipal profiles $\bm{\phi}_s=(\bm{\beta}_s,\xi_s)$ are clustered through
\(\bm{\phi}_s\mid G\sim G,
G\sim\mathrm{DP}(\alpha_{\mathrm{DP}},G_0),
\)
where $\alpha_{\mathrm{DP}}\sim\mathrm{Gamma}(3,2)$ controls cluster concentration and
$G_0=\mathcal{N}(\mathbf{0},I)\times U(-1,1)$ is the base distribution, allowing data-driven groups without administrative boundaries.

\item \textbf{Administrative hierarchical baseline (\model{M6}):}\\
\noindent IID region and province effects plus a national AR(1) temporal effect,
without using the municipal adjacency graph $W$.

\end{itemize}
\end{block}

\vspace{0.5cm}

%------------------------------------------------------------------------------
\begin{block}{6. Out-of-sample model comparison}

\justifying
Two rolling-origin \textbf{one-year-ahead evaluations}:
$2004$--$2017\!\rightarrow\!2018$ and
$2004$--$2018\!\rightarrow\!2019$
($N_{\mathrm{test}}=15{,}792$).

\vspace{0.5cm}
\centering
\footnotesize
\setlength{\tabcolsep}{3.5pt}
\begin{tabular}{lcccccc}
\toprule
\textbf{Model} & \textbf{RMSE} & \textbf{MAE} &
\textbf{Cov$_{95}$} & \textbf{IS$_{95}$} &
\textbf{Time} & \textbf{MaxRSS} \\
\midrule
\model{M1} INLA Global   & \textbf{13.69} & \textbf{8.81} & 0.961 & 84.65 & 00:18:54 & 107.9 GiB \\
\model{M2} MCMC Global   & 13.69 & 8.82 & 0.961 & 84.63 & 05:48:07 & 2.6 GiB \\
\model{M3} Overlap INLA  & 13.96 & 9.09 & 0.962 & \textbf{77.09} & 107 jobs & -- \\
\model{M4} Disjoint INLA & 14.20 & 9.34 & 0.962 & 77.60 & 107 jobs & -- \\
\model{M5} BSTC-DP       & 17.18 & 10.83 & 0.931 & 99.30 & 06:40:10 & 4.2 GiB \\
\model{M6} Admin IID     & 13.77 & 8.90 & \textbf{0.972} & 88.33 & 00:01:10 & 1.9 GiB \\
\bottomrule
\end{tabular}

{\scriptsize\textit{
Predictive metrics pooled over both folds; computational costs refer to the 2019 fold.}}

\vspace{0.5cm}
\justifying
\normalsize
\begin{itemize}
  \item \model{M1} and \model{M2} give nearly identical point accuracy, with the lowest RMSE and MAE;
  \item \model{M3}/\model{M4} achieve the lowest IS$_{95}$, indicating the
  best sharpness--coverage trade-off;
  \item \model{M5} shows poorer out-of-sample accuracy and a highly fragmented posterior partition ($K\approx130$);
  \item \model{M6} reaches the highest coverage, but with a higher IS$_{95}$.
\end{itemize}
\end{block}

\smallgap
%------------------------------------------------------------------------------
\begin{block}{Conclusions}

\begin{itemize}

  %\item \textbf{Substantive:}
  %Municipal net migration shows interpretable \textbf{covariate associations} together with substantial \textbf{latent spatial dependence} ($\rho\approx0.77$) and \textbf{temporal persistence} ($\xi\approx0.61$).

  %\item \textbf{Methodological:}
  %The global CAR$\times$AR(1) model gives the best point prediction, whereas the partitioned models provide sharper predictive intervals and reveal local heterogeneity.


\item \textbf{Covariate associations:}
the strongest associations with net migration are positive for employment and population density, and negative for accessibility deficits and unemployment.

\item \textbf{Spatio-temporal structure:}
persistent spatial patterns emerge, with substantial spatial dependence ($\rho \approx 0.77$) and temporal persistence ($\xi \approx 0.61$).
\item \textbf{Territorial heterogeneity:}
covariate effects vary across provinces in both magnitude and, for some predictors, sign.


\end{itemize}

\end{block}
\end{column}

%==============================================================================
% COLUMN 3: CHOSEN MODEL (M1), DIAGNOSTICS & POLICY IMPLICATIONS
%==============================================================================
\begin{column}{0.324\textwidth}

%------------------------------------------------------------------------------
\begin{block}{7. Associations with net migration}

\justifying
M1/M2 estimate common national covariate associations, while municipality-level residual variation is captured by a spatio-temporal field with common dependence parameters. %It achieves the lowest pooled RMSE and MAE.

%The global model (\model{M1}/\model{M2}) is used as the \textbf{national reference}, attaining the lowest pooled RMSE and MAE while retaining the full municipal spatio-temporal structure.

\microgap
\begin{center}
  \includegraphics[width=0.98\linewidth]
  {Poster_figures/m1_rw1_coefficients_intervals.png}
  \vspace{-0.1cm}
  
  \parbox{0.96\linewidth}{
  \footnotesize
  \textit{Figure 2: Posterior 95\% credible intervals for \model{M1} fixed effects.}
  }
\end{center}

\microgap
\justifying
\textbf{Substantive interpretation of the fixed effects:}

\begin{itemize}
    \item \textbf{Attraction factors:}
    employment rate shows the strongest positive association with net migrationb($\hat{\beta}=+0.250$), followed by mean age ($\hat{\beta}=+0.074$).

    \item \textbf{Repulsion factors:}
    service inaccessibility ($\hat{\beta}=-0.053$), changes in the foreign-resident share
    ($\hat{\beta}=-0.042$), and unemployment ($\hat{\beta}=-0.026$)
    are associated with lower net migration.

    \item \textbf{Limited evidence of association:}
unemployment and GDP growth show 95\% credible intervals overlapping zero.
\end{itemize}

\end{block}

\smallgap
\begin{block}{8. Spatio-temporal structure and residual diagnostics}

\begin{center}
  \includegraphics[width=0.98\linewidth]
  {Poster_figures/m1_latent_field_municipal.png}
  \vspace{-0.1cm}

  \parbox{0.96\linewidth}{
  \footnotesize
  \textit{Figure 3: Posterior mean of the latent field $\mathbf{w}_t$ at municipal level, by year (M1/M2).}
  }
\end{center}

\vspace{0.3cm}
\centering
\footnotesize
\setlength{\tabcolsep}{3.5pt}
\renewcommand{\arraystretch}{0.92}

\begin{tabular}{lcccc}
\toprule
\textbf{Parameter} & \textbf{Mean} & \textbf{SD} & \textbf{2.5\%} & \textbf{97.5\%} \\
\midrule
Spatial dependence $\rho$ & \textbf{0.7774} & 0.0263 & 0.7197  & 0.8224  \\
Temporal persistence $\xi$ & \textbf{0.6118} & 0.0144 & 0.5825  & 0.6392  \\
Latent field variance $\tau^2_{ST}$ & 0.3623          & 0.0268 & 0.3139  & 0.4190  \\
Residual variance $\sigma^2_\epsilon$ & 0.6689          & 0.0086 & 0.6521  & 0.6860  \\
Random walk variance $\tau^2_{RW1}$ & 0.1808         & 0.0809 & 0.0562 & 0.3636 \\
\bottomrule

\end{tabular}

{\scriptsize\textit{
Common spatio-temporal parameters (M1/M2).}}

\renewcommand{\arraystretch}{1}
\microgap
\justifying
\textbf{Latent spatial structure and residual diagnostics:}

\begin{itemize}

\item \textbf{Persistent spatial pattern:}
the latent field shows marked territorial clustering, with spatial patterns
that persist but evolve over time.

\item \textbf{Spatial residual adequacy:}
residual Moran's $I$ remains close to zero and slightly negative
($-0.07$ to $-0.03$), indicating no remaining positive spatial autocorrelation.

\item \textbf{Temporal residual stability:}
annual mean residuals remain close to zero throughout the study period,
with no marked residual temporal trend.

\end{itemize}

\end{block}

%------------------------------------------------------------------------------

\begin{block}{9. Local spatial heterogeneity}
\begin{center}
  \includegraphics[width=0.98\linewidth]
  {Poster_figures/M3_selected_local_effects_5x1.png}
  \vspace{-0.15cm}

  \parbox{0.96\linewidth}{
  \scriptsize
  \textit{Figure 4: Posterior means of selected province-specific coefficients
 from the M3 overlap fit.}
  }
\end{center}

\microgap
\justifying
\textbf{Beyond the global coefficients:}
\begin{itemize}

  \item \textbf{Employment} remains predominantly positively associated with
  net migration, but its magnitude varies markedly across provinces.

  \item Several associations \textbf{vary locally in both magnitude and sign}, revealing heterogeneity beyond a single national coefficient.

\end{itemize}

\end{block}

\smallgap
%------------------------------------------------------------------------------
\begin{block}{What's next?}
\justifying
\begin{itemize}

\item \textbf{Variable selection:}
learn relevant covariates rather than prespecifying them;

\item \textbf{Robust tails modeling:}
quantile or heavy-tailed space--time models to better capture extreme migration dynamics;

\item \textbf{DP clustering refinement:}
use heavy-tailed base distributions or Gaussian-mixture structure learned from
partition-specific coefficients to inform the clustering prior.

\end{itemize}

\end{block}
\end{column}
\end{columns}

%==============================================================================
% FOOTER
%==============================================================================

\vfill
\hrule height 1pt
\vspace{0.20cm}

\begin{center}
\footnotesize
\textbf{Contacts:}  claudia.franceschini@phd.unipd.it \quad$\cdot$\quad sophiegrace.parolin@phd.unipd.it \quad$\cdot$\quad teresa.secco@unive.it
\qquad\qquad
\textbf{Department of Statistical Sciences, University of Padua}
\end{center}

\end{frame}
\end{document}
